% ==========================================================================
% Capítulo: Marco teórico
% Estado: se registra en ../STATUS.md, no aquí.
%
% Debe contener (project-context/formato-anteproyecto.md, "Marco teórico"):
%   - Fundamentos, conceptos teóricos y técnicos propios de la disciplina
%     sobre los que se apoya el proyecto.
%   - Marcos de referencia y normativos vigentes cuando apliquen.
%   - OBLIGATORIO: explicitar cómo se relaciona CADA concepto definido con
%     los objetivos específicos y/o la metodología definida.
% NO debe contener: cadena de resúmenes de artículos (eso ni siquiera es
% estado del arte), ni conceptos que no se usen después en el documento.
%
% EXTENSIÓN (decisión de los autores, 2026-09-13): el formato fija 4 páginas
% como máximo. Los autores levantaron ese tope para este capítulo y fijaron
% uno propio de 8 PÁGINAS. No lo recorte a 4 ni lo deje crecer más allá de 8
% sin consultarlo con ellos.
%
% PESOS RELATIVOS (decisión de los autores, 2026-09-13): los recursos
% personalizados y el patrón operador se explican, pero NO dominan el
% capítulo; la identidad federada y el control de acceso ocupan dos párrafos;
% contenedores y orquestación va resumido. Si hay que recortar de nuevo,
% recorte de ahí antes que del ciclo de vida o del servicio de inferencia.
%
% ORGANIZACIÓN: bases teóricas continuas, de lo general a lo específico
% (ciclo de vida de IA/ML -> despliegue e inferencia -> contenedores y
% orquestación -> extensibilidad declarativa -> gobernanza -> identidad ->
% observabilidad), y cierre con MARCO CONCEPTUAL Y CATEGORÍAS DE ANÁLISIS,
% que concentra en una tabla el vínculo concepto -> objetivo específico que
% exige el formato. NO cierre cada párrafo con un \cref{obj:...}: eso hacía
% ilegible la primera versión.
% Los ANTECEDENTES de investigación NO van aquí: viven en el cap. 01
% (contexto institucional) y en el cap. 06 (estado del arte).
%
% Fuentes: IEEE, ACM, libros acreditados. Citar con natbib
% (../../skills/latex/references/citations-and-figures.md); toda afirmación traza
% a una fuente real (../../agent-roles/Investigador.md).
% ==========================================================================
\chapter{Marco teórico}\label{cha:marco-teorico}

Construir una plataforma que despliegue modelos de inteligencia artificial sobre hardware compartido, gobierne quién los ejecuta y observe lo que ocurre en la infraestructura mientras lo hacen exige integrar conocimiento de cinco dominios de la ingeniería de sistemas: el ciclo de vida de los sistemas de aprendizaje automático, el servicio de modelos en producción, la orquestación de contenedores y su modelo de extensión, la gobernanza de recursos compartidos entre usuarios con derechos desiguales, y la observabilidad de infraestructura acelerada. Este capítulo los recorre de lo general a lo particular, desde la fase del ciclo de vida que el proyecto ocupa hasta las técnicas concretas que la instrumentan, y cierra con las definiciones operativas y las categorías bajo las cuales se analiza el resultado.

\section{El ciclo de vida de las cargas de inteligencia artificial}\label{sec:mlops-ciclo-vida}

MLOps designa la disciplina que extiende al ciclo de vida completo de un sistema de aprendizaje automático las prácticas de integración y entrega continuas, control de versiones y monitoreo propias del desarrollo de software convencional \citep{kreuzberger-mlopsoverview-2023}, y ubica el despliegue, el monitoreo en producción y la gobernanza de acceso entre sus prácticas centrales \citep{lima-mlopspractices-2022}. La razón por la que estas prácticas concentran tanta atención la documentaron \citet{sculley-hiddentechnicaldebt-2015}: en un sistema de aprendizaje automático real, el código del modelo ocupa una fracción pequeña del conjunto, y todo lo que lo rodea, la gestión de configuración, el servido, la infraestructura de cómputo y el monitoreo, constituye el grueso del sistema. Tratar el modelo entrenado como si fuera el producto terminado genera una deuda técnica que permanece invisible durante el desarrollo y emerge cuando el sistema debe operar de forma sostenida.

Dentro de ese ciclo, la literatura separa dos fases cuyos objetivos de diseño se oponen. El entrenamiento consume datos para producir un artefacto de pesos, maximiza el throughput sobre un horizonte largo y no interactivo, tolera la interrupción y se mide por el tiempo hasta converger. El servicio de modelos, o \emph{model serving}, parte de un artefacto ya producido, atiende solicitudes individuales en tiempo real y se mide por latencia y por estabilidad bajo carga variable \citep{muiruri-mlinferenceserving-2026}. La distinción es general y no pertenece al dominio de los modelos de lenguaje: un clasificador clásico, un modelo de visión por computador y un modelo generativo de gran escala atraviesan el mismo ciclo y presentan la misma separación entre la fase que produce el artefacto y la fase que lo sirve. Lo que cambia entre ellos es el perfil de consumo y el conjunto de métricas relevantes, sin que la estructura del ciclo se altere. El proyecto se ubica por completo en la fase de servicio, para cargas de inteligencia artificial en general, mientras el entrenamiento permanece bajo el alcance del proyecto hermano del mismo macro-proyecto.

Esa frontera conceptual necesita un criterio operativo, porque ambas fases se ejecutan sobre las mismas máquinas de la sala 104M y una plataforma que gobierne el hardware no puede quedar ciega ante la mitad de lo que corre en él. El proyecto adopta un criterio de punto de envío: un trabajo cuenta como carga de entrenamiento cuando ingresa al clúster por el plano de control del proyecto hermano, y como carga de despliegue cuando ingresa por el plano de control de esta plataforma. La clase queda determinada por la forma en que el trabajo entró, sin depender de un umbral observado durante su ejecución. La alternativa, inferirla de la telemetría mediante una heurística del tipo memoria de video sostenida sin tráfico de inferencia, produce clasificaciones que cambian con la carga del momento y que resultan imposibles de auditar después.

\section{Despliegue e inferencia de modelos}\label{sec:despliegue-inferencia}

Desplegar un modelo consiste en instalar un artefacto ya entrenado en un entorno de ejecución donde pueda recibir solicitudes de forma sostenida, e inferir es la operación por la cual ese artefacto produce una salida a partir de una entrada nueva. Entre ambos se interpone el motor de inferencia: el componente que carga los pesos en la memoria del acelerador, gestiona el conjunto de solicitudes concurrentes y expone una interfaz de red por la que los clientes consumen el modelo. \citet{muiruri-mlinferenceserving-2026} lo caracterizan como el punto donde se concentran las decisiones de rendimiento del sistema en producción, porque condiciona la latencia percibida, la concurrencia que un nodo sostiene y la memoria que un modelo ocupa realmente. Las métricas con que se evalúa el servicio se derivan de esa caracterización: el tiempo hasta el primer token, la latencia entre tokens, el throughput agregado y el tiempo de espera en cola. Las cuatro guardan dependencias entre sí, porque aumentar el tamaño del lote eleva el throughput agregado y a la vez degrada la latencia entre tokens de cada solicitud, de modo que el motor opera siempre sobre un compromiso explícito.

Cuando el modelo servido es un modelo de lenguaje de gran escala, el motor enfrenta un cuello de botella particular: la caché de atención crece con la longitud de cada petición y, bajo solicitudes concurrentes de longitud variable, fragmenta la memoria del acelerador hasta desperdiciar una fracción sustancial de la capacidad instalada. \citet{kwon-pagedattention-2023} resuelven esa fragmentación paginando la caché en bloques de tamaño fijo, con una analogía directa a la memoria virtual de un sistema operativo. \citet{yu-orca-2022} atacan el mismo problema desde la programación de peticiones y muestran que planificar a nivel de iteración, en lugar de esperar a que el lote completo termine, eleva de forma sustancial el throughput agregado del servidor. \citet{miao-llmservingsurvey-2025} organizan estas técnicas, junto con la compresión de modelos y la programación de kernels, como un espectro continuo de optimización que abarca desde el algoritmo hasta el sistema.

Dentro de ese espectro, la cuantización posentrenamiento ocupa un lugar particular para este proyecto. La técnica reduce la precisión numérica con la que se representan los pesos, con una degradación de calidad acotada y medida experimentalmente (\citealp{frantar-gptq-2023}; \citealp{lin-awq-2024}; \citealp{dettmers-llmint8-2022}; \citealp{chen-llmquantizationsurvey-2026}). Su relevancia aquí es de capacidad y no de construcción: explica por qué modelos de 7\,000 a 14\,000 millones de parámetros caben dentro de los 24\,GB de memoria de video de un nodo con tarjeta RTX 4090, y por tanto por qué el laboratorio puede servirlos sin recurrir a inferencia distribuida entre nodos. La plataforma despliega pesos que ya llegan cuantizados, o se apoya en la cuantización que el propio motor aplica al cargar el modelo, como hace \emph{llama.cpp}, sin ejecutar ningún proceso propio de cuantización.

Los motores disponibles ocupan puntos distintos de ese espectro, y esa diversidad es la razón teórica por la cual la plataforma no puede acoplarse a ninguno. vLLM aplica paginación de caché y programación por iteración para maximizar el throughput bajo concurrencia alta, y paga ese throughput con mayor latencia entre tokens cuando el lote crece \citep{kwon-pagedattention-2023, miao-llmservingsurvey-2025}. \emph{llama.cpp} adopta el compromiso contrario y prioriza la latencia de un único flujo junto con una huella de despliegue mínima, lo que lo hace adecuado para sesiones individuales sobre hardware de clase estación de trabajo. Otros motores sirven familias de modelos que no son de lenguaje, con formatos de artefacto y perfiles de cómputo propios. La literatura revisada no ofrece una comparación controlada entre estas opciones sobre el tipo de hardware que el IAsLab tiene instalado, de modo que cualquier preferencia operativa del laboratorio constituye una decisión de ingeniería propia y no un resultado establecido. Esa ausencia obliga al proyecto a tratar el motor como pieza intercambiable y a producir la comparación mediante medición.

\section{Contenedores y orquestación de cargas de cómputo}\label{sec:orquestacion-contenedores}

Un motor de inferencia, con sus pesos, sus dependencias de CUDA y sus versiones concretas de bibliotecas, constituye una unidad de software cuya reproducción por instalación manual es frágil. El contenedor la empaqueta junto con su entorno de ejecución completo en una imagen inmutable, de modo que lo que se ejecuta en un nodo es idéntico a lo que se probó en otro; el orquestador resuelve el problema siguiente, que consiste en decidir en qué máquina se ejecuta cada unidad, reiniciarla cuando falla y exponerla a la red de forma estable. Kubernetes se consolidó como el orquestador de referencia para esa tarea: un análisis bibliométrico identifica la programación de recursos, el cómputo distribuido y la inteligencia artificial entre sus focos de investigación más activos \citep{carrion-k8sbibliometric-2022}, y una revisión multivocal documenta, junto a los beneficios que los profesionales reportan, la carencia de herramientas nativas de diagnóstico entre sus obstáculos recurrentes \citep{shamim-k8smultivocal-2022}. Ese balance sostiene dos decisiones del proyecto a la vez: adoptar Kubernetes como sustrato en lugar de construir un programador de clúster propio, y prever desde el diseño la capa de observabilidad que cubra el vacío de diagnóstico que la propia literatura señala.

Los primitivos que el orquestador ofrece se aplican de forma directa al caso del laboratorio. Los selectores de nodo, junto con las marcas y tolerancias, dirigen las cargas que requieren aceleración a las máquinas con tarjeta gráfica y reservan las demás para los servicios de control. Los espacios de nombres dividen el clúster en fronteras lógicas de administración sobre las que se aplican límites agregados de consumo. Los volúmenes persistentes desacoplan el almacenamiento de los pesos del ciclo de vida de los contenedores que los sirven, lo que evita descargar decenas de gigabytes cada vez que un despliegue se reinicia. Ninguno de estos mecanismos es específico de cargas de inteligencia artificial, y esa generalidad es lo que permite gobernar por igual cargas de inferencia, cargas de entrenamiento del proyecto hermano y servicios auxiliares.

\section{Extensibilidad declarativa del orquestador}\label{sec:crd-operadores}

Kubernetes no fue diseñado para cargas de inteligencia artificial y sin embargo las hospeda sin que su núcleo haya sido rediseñado. El mecanismo que lo permite es su modelo declarativo. \citet{burns-borgomegak8s-2016}, a partir de una década operando los sistemas Borg, Omega y Kubernetes, describen su principio central: el usuario declara el estado que desea que el clúster tenga, y un controlador ejecuta de forma permanente un bucle de reconciliación que compara ese estado con el observado y actúa hasta hacerlos converger. Como toda acción se decide a partir de la observación del estado presente, y no de una máquina de estados que recuerde por dónde iba el proceso, un controlador que falla y se reinicia retoma el trabajo desde lo que encuentra. Esa propiedad es la que permite operar una infraestructura descrita íntegramente en archivos de configuración versionados y redesplegarla sin reconstruirla a mano, algo que el proyecto necesita porque el clúster se levanta desde cero sobre las estaciones de trabajo de la sala.

Sobre ese modelo se apoya la vía de extensión. Una definición de recurso personalizado (\emph{Custom Resource Definition}, CRD) registra en el servidor de API un tipo de objeto que Kubernetes no conocía, junto con el esquema que valida sus campos, de modo que un concepto propio de un dominio pasa a ser un objeto de primera clase del clúster. El controlador que acompaña a ese tipo de recurso, y que encapsula el conocimiento operativo necesario para materializarlo, recibe el nombre de operador. \citet{xu-k8soperatorbugs-2024} estudiaron 210 defectos recolectados de 36 operadores de uso extendido y encontraron que sus fallas más frecuentes provienen del manejo incorrecto del estado durante la reconciliación. De ese hallazgo se sigue una decisión de ingeniería explícita para el proyecto: componer operadores ya maduros para todo aquello que alguno de ellos cubra, y reservar el desarrollo propio para la lógica que ninguno provee.

Las extensiones que el proyecto compone cubren tres necesidades. El NVIDIA GPU Operator aprovisiona de forma declarativa el controlador de dispositivo, el runtime de contenedores y el complemento que expone la tarjeta gráfica al programador del clúster como recurso asignable, trabajo que de otro modo exigiría una instalación manual por nodo. KubeRay, sobre el framework de cómputo distribuido que \citet{moritz-ray-2018} describen a partir de tareas asíncronas y actores con estado, aporta los recursos con los que un clúster de cómputo, un trabajo por lotes y un servicio de inferencia se declaran como objetos nativos, de modo que desplegar un modelo deja de ser la ejecución de un script sobre una máquina concreta. Una capa de admisión y encolamiento, por último, retiene los trabajos hasta que el clúster puede garantizarles la capacidad completa que solicitan, lo que evita que varias cargas parcialmente aprovisionadas se bloqueen entre sí; \citet{gao-dlschedulingtaxonomy-2022} sitúan ese mecanismo entre las políticas de admisión y expropiación propias de los clústeres de aprendizaje profundo, y es además el punto donde se materializan las políticas de gobernanza de la \cref{sec:gobernanza-recursos-compartidos}.
% Candidato concreto para esa capa: Kueue, señalado por los autores en
% busquedas-realizadas.md. NO se nombra en la prosa porque todavía no está
% en project-context/technologies.md; nombrarlo aquí lo convertiría en un
% compromiso arquitectónico que nadie ha tomado (CLAUDE.md regla 6).

\section{Gobernanza de recursos compartidos}\label{sec:gobernanza-recursos-compartidos}

Una infraestructura de cómputo acelerado compartida entre usuarios con necesidades desiguales plantea un problema de asignación que la literatura de sistemas ha abordado desde dos familias con supuestos opuestos. La primera reparte la capacidad de forma proporcional entre solicitantes considerados equivalentes y optimiza alguna métrica agregada del clúster: \citet{xiao-gandiva-2018} explotan la estructura iterativa de las cargas de aprendizaje profundo para migrar y reempaquetar trabajos en ejecución, \citet{gu-tiresias-2019} asignan sin conocer de antemano la duración de cada trabajo, y \citet{mahajan-themis-2020}, \citet{narayanan-gavel-2020} y \citet{qiao-pollux-2021} formulan la asignación como un problema de equidad a largo plazo, generalizado a aceleradores heterogéneos y acoplado a los parámetros de cada entrenamiento, sobre la base de sistemas de gestión de clústeres como el que documentan \citet{verma-borg-2015}. Lo que todas comparten es el supuesto de que los solicitantes tienen el mismo derecho sobre el recurso.

La segunda familia parte del supuesto contrario: establece jerarquías explícitas entre clases de solicitante y garantiza acceso preferente a las superiores mediante prioridad y, cuando hace falta, expropiación de recursos ya asignados. \citet{gao-dlschedulingtaxonomy-2022} tratan ambas familias por separado en su taxonomía porque persiguen objetivos incompatibles, equidad entre pares en un caso y garantía de acceso en el otro, de modo que la elección entre ellas responde a una definición organizacional antes que técnica. El contexto académico del IAsLab corresponde sin ambigüedad a la segunda: estudiantes y profesores reservan capacidad, pero los profesores requieren prioridad para reclamar recursos y para controlar el reparto de la capacidad de la sala, incluso por encima de una reserva estudiantil activa, porque la docencia y la investigación dirigida tienen precedencia institucional sobre la exploración individual. Ninguna política de reparto proporcional entre iguales satisface ese requisito, por bien que optimice la utilización agregada.

Dos mecanismos instrumentan la gobernanza que el proyecto necesita. La reserva anticipada aparta capacidad por adelantado para una carga específica durante una ventana de tiempo definida; \citet{weitzel-nrpreservation-2025} documentan su introducción sobre unidades de procesamiento gráfico en un clúster Kubernetes multiusuario de escala nacional y reportan un aumento sustancial de la utilización promedio, lo que muestra que la reserva mejora el aprovechamiento del hardware al reducir el tiempo que los recursos pasan ociosos entre usos no coordinados. El sobreaprovisionamiento controlado la complementa al asignar más capacidad nominal de la instalada, bajo el supuesto estadístico de que los solicitantes no agotan su cuota máxima de forma simultánea; el problema se ha formalizado como un empaquetado bajo restricciones probabilísticas \citep{cohen-cloudovercommit-2019} y desplegado a escala de producción en múltiples centros de datos \citep{bashir-peakovercommit-2021}, y su viabilidad descansa sobre la predictibilidad del consumo real frente al declarado, que \citet{cortez-resourcecentral-2017} establecen al caracterizar las cargas de una plataforma de nube a gran escala. El margen que el proyecto adopta, entre el 10\,\% y el 20\,\%, proviene de la experiencia operativa del laboratorio y no de ninguna de estas fuentes.

En el plano de la implementación, las cuotas nativas que Kubernetes asocia a los espacios de nombres acotan el consumo agregado de procesador, memoria y dispositivos de un conjunto de cargas, y constituyen el punto de aplicación sobre el que se construyen la reserva y el margen de sobreasignación: la cuota define cuánto puede consumir un espacio de nombres, y la capa de admisión de la \cref{sec:crd-operadores} decide en qué momento un trabajo concreto obtiene ese consumo. Conviene señalar, por último, que gobernar cómputo acelerado en una institución educativa tiene rasgos que la literatura de centros de datos comerciales no captura. \citet{george-hpcgpuclassroom-2020} revisan el uso de unidades de procesamiento gráfico en contextos de aula e investigación y sitúan el problema en un marco donde el objetivo consiste en sostener el acceso de una población estudiantil rotativa, con demandas intermitentes y desiguales, antes que en maximizar el retorno económico por hora de acelerador. Ese encuadre es el que corresponde al IAsLab y el que justifica organizar la gobernanza alrededor de reservas y prioridades de rol.

\section{Identidad federada y control de acceso}\label{sec:identidad-rbac}

Aplicar una política de prioridad por rol exige saber con certeza quién solicita un recurso y qué rol ostenta. El control de acceso basado en roles sigue siendo el modelo de referencia para autorización en sistemas multiusuario, con extensiones para cómputo en la nube que incorporan atributos del solicitante y niveles de confianza a la decisión \citep{li-rbaccloudsurvey-2015}. La autenticación federada resuelve el problema complementario al delegar la verificación de la identidad en un proveedor externo, que tras autenticar al usuario emite un token firmado con sus atributos, entre ellos sus roles; \citet{fett-oauth2security-2016} establecen mediante análisis formal las garantías de autorización, autenticación e integridad de sesión que ofrece OAuth 2.0 bajo ese esquema. La Universidad Icesi opera un proveedor institucional, SAAMFI, comparable en función a productos como Keycloak o Auth0. Una plataforma apoyada en él hereda su taxonomía de roles y define por encima solo los que su propia operación requiere, como el de administrador, lo que evita comprometerse con una enumeración de perfiles académicos que pertenece al proveedor y puede cambiar sin intervención del proyecto.

La aplicación de cuotas sobre el tráfico de inferencia se concentra, además, en un único punto de entrada, según el patrón que la literatura de microservicios documenta como \emph{API gateway}: un componente que agrega las interfaces de los servicios subyacentes, aplica sobre el tráfico las políticas transversales de autenticación, autorización y límite de tasa, y registra el consumo \citep{montesi-apigatewaypatterns-2016}. Su variante especializada para modelos de lenguaje añade límites basados en el consumo de tokens, que es la unidad en la que se contabiliza el uso de ese tipo de modelos. La revisión de esta literatura no localizó ninguna variante que aplique cuotas de forma unificada sobre cargas de lenguaje y cargas que no lo son, vacío relevante para el proyecto porque la plataforma debe gobernar cualquier carga de inteligencia artificial, y un mecanismo que solo sepa contar tokens deja fuera de la política a todo modelo cuya salida no se mida en esa unidad.

\section{Observabilidad y telemetría de infraestructura acelerada}\label{sec:observabilidad-telemetria}

La observabilidad de un sistema distribuido se define como la capacidad de explicar su comportamiento interno a partir de lo que puede medirse desde afuera, y se apoya en tres señales complementarias: las métricas, que agregan magnitudes numéricas a lo largo del tiempo; las trazas, que siguen el recorrido de una solicitud a través de los componentes que la atienden; y los registros, que conservan los eventos discretos que cada componente emite \citep{li-observabilitysurvey-2021}. Las tres resultan necesarias porque responden preguntas distintas: una métrica indica que algo se degradó, una traza indica dónde, y un registro indica qué ocurrió en ese punto.

En un clúster de cómputo acelerado compartido, esa capacidad depende de telemetría de hardware de grano fino, y la evidencia empírica es abundante. \citet{weng-mlaasinthewild-2022} documentan, en un clúster de producción a gran escala, patrones sostenidos de subutilización que solo la instrumentación por nodo permite detectar porque las métricas agregadas los ocultan; \citet{jeon-philly-2019} muestran que una parte sustancial del tiempo total de los trabajos se consume en espera y en fallos, visible únicamente cuando se instrumentan las colas además de los nodos; \citet{gao-lowgpuutilization-2024} encuentran que la mayoría de los casos de baja utilización obedece a causas corregibles identificables a partir de la telemetría; y \citet{liu-gpufailureprediction-2022} muestran que un volumen extenso de telemetría recolectada durante meses eleva de forma sustancial la precisión con que se anticipan los fallos de hardware. De este conjunto se desprenden las señales primarias para diagnosticar saturación y anticipar fallo: temperatura del núcleo y de la memoria, utilización de memoria de video, consumo de potencia y throughput del bus PCIe.

El ecosistema de Kubernetes instrumenta esas señales con un conjunto de componentes establecido. Prometheus recolecta métricas por consulta periódica a los puntos de exposición de cada componente, las almacena como series temporales identificadas por etiquetas y evalúa sobre ellas reglas de alerta. Loki agrega los registros que emiten los contenedores y los indexa por las mismas etiquetas, lo que permite pasar de una anomalía observada en una gráfica a los registros que la produjeron sin cambiar de sistema de coordenadas. Grafana consolida ambas fuentes en paneles de visualización. El exportador DCGM de NVIDIA completa el arreglo al exponer a Prometheus las métricas de bajo nivel de las tarjetas aceleradoras, que residen en el firmware del dispositivo y ningún otro componente puede obtener. Una condición de diseño atraviesa el conjunto y conecta de vuelta con la \cref{sec:mlops-ciclo-vida}: la telemetría de un nodo no puede presuponer que la carga activa es de inferencia, porque ese mismo nodo puede estar saturado por un trabajo de entrenamiento sometido a través de otro plano de control, y un panel que asuma lo contrario atribuirá el consumo al usuario equivocado. Rotular cada métrica con la clase de carga que la originó, según el criterio de punto de envío ya definido, permite separar ambos consumos dentro de un mismo panel sin inferir nada, y hace auditable la aplicación de las cuotas sobre ambas clases.

\section{Marco conceptual y categorías de análisis}\label{sec:marco-conceptual}

Los conceptos desarrollados en las secciones anteriores se emplean en el resto del documento con el significado operativo acotado que fija la \cref{tab:marco-conceptual}, junto con el objetivo específico que cada uno sostiene. La tabla cumple el requisito del formato de relacionar cada concepto definido con los objetivos específicos, y lo concentra en un solo lugar para que la fundamentación de las secciones previas pueda leerse como un argumento continuo.

{\small
\begin{xltabular}{\textwidth}{p{0.18\textwidth} X p{0.13\textwidth}}
        \caption{Definición operativa de los conceptos del marco teórico y objetivo específico que los emplea.}
        \label{tab:marco-conceptual} \\
        \toprule
        \textbf{Concepto} & \textbf{Definición operativa en este proyecto} & \textbf{Objetivo} \\
        \midrule
        \endfirsthead
        \toprule
        \textbf{Concepto} & \textbf{Definición operativa en este proyecto} & \textbf{Objetivo} \\
        \midrule
        \endhead
        \bottomrule
        \endfoot
        Carga de despliegue &
        Trabajo que ingresa al clúster por el plano de control de esta plataforma y sirve un modelo ya entrenado para atender solicitudes de inferencia &
        \Cref{obj:orquestacion} \\
        \addlinespace
        Carga de entrenamiento &
        Trabajo que ingresa por el plano de control del proyecto hermano. El proyecto no lo implementa, pero lo gobierna y lo observa en igualdad de condiciones &
        \Cref{obj:telemetria}, \cref{obj:gobernanza} \\
        \addlinespace
        Motor de inferencia &
        Componente que carga los pesos de un modelo, gestiona la concurrencia y expone una interfaz de red, empaquetado como imagen de contenedor e intercambiable sin rediseñar la plataforma &
        \Cref{obj:orquestacion} \\
        \addlinespace
        Recurso personalizado y operador &
        Tipo de objeto registrado en el servidor de API y controlador que lo reconcilia. Es el mecanismo por el que la plataforma declara clústeres de cómputo, trabajos y servicios de inferencia como objetos nativos del orquestador &
        \Cref{obj:orquestacion} \\
        \addlinespace
        Admisión y encolamiento &
        Retención de un trabajo hasta que el clúster puede garantizarle la capacidad completa que solicita. Es el punto donde se materializa la política de prioridad &
        \Cref{obj:gobernanza} \\
        \addlinespace
        Reserva con prioridad de rol &
        Asignación anticipada de capacidad por ventana de tiempo, con un margen de sobreasignación del 10\,\% al 20\,\% y con potestad de los profesores para reclamar recursos y controlar el reparto de la capacidad de la sala &
        \Cref{obj:gobernanza} \\
        \addlinespace
        Identidad federada &
        Delegación de la autenticación en SAAMFI, cuyos roles la plataforma hereda y sobre los cuales define un único rol adicional de administración &
        \Cref{obj:gobernanza} \\
        \addlinespace
        Telemetría agnóstica a la carga &
        Captura de métricas de hardware y de trabajo por nodo, rotuladas con la clase de carga que las originó, con independencia de si el nodo ejecuta inferencia o entrenamiento &
        \Cref{obj:telemetria} \\
        \addlinespace
        Aceptación operacional &
        Grado en que la plataforma sostiene la concurrencia objetivo de veinte usuarios simultáneos y alcanza un nivel satisfactorio de usabilidad percibida &
        \Cref{obj:evaluacion} \\
\end{xltabular}
}

Las dos últimas filas fijan las categorías bajo las cuales se analiza el resultado del proyecto, cada una con un instrumento de respaldo metodológico propio. El desempeño se observa a través de las métricas de hardware e inferencia que la propia plataforma produce, recogidas bajo pruebas de carga cuyo diseño, ejecución y análisis siguen la sistematización que \citet{jiang-loadtestingsurvey-2015} elaboran sobre la literatura de pruebas de carga en sistemas de gran escala. La aceptación operacional se observa mediante el \emph{System Usability Scale}, cuestionario de diez ítems cuyas bandas de referencia por ítem recalibraron \citet{lewis-susbenchmarks-2018} a partir de 446 estudios y más de cinco mil respuestas, lo que permite interpretar el puntaje frente a una línea base publicada en lugar de valorarlo de forma aislada. El \cref{cha:metodologia} desarrolla la aplicación concreta de ambos dentro de la fase de evaluación empírica.
